% Options for packages loaded elsewhere
\PassOptionsToPackage{unicode}{hyperref}
\PassOptionsToPackage{hyphens}{url}
\documentclass[
]{article}
\usepackage{xcolor}
\usepackage[margin=1in]{geometry}
\usepackage{amsmath,amssymb}
\setcounter{secnumdepth}{-\maxdimen} % remove section numbering
\usepackage{iftex}
\ifPDFTeX
  \usepackage[T1]{fontenc}
  \usepackage[utf8]{inputenc}
  \usepackage{textcomp} % provide euro and other symbols
\else % if luatex or xetex
  \usepackage{unicode-math} % this also loads fontspec
  \defaultfontfeatures{Scale=MatchLowercase}
  \defaultfontfeatures[\rmfamily]{Ligatures=TeX,Scale=1}
\fi
\usepackage{lmodern}
\ifPDFTeX\else
  % xetex/luatex font selection
\fi
% Use upquote if available, for straight quotes in verbatim environments
\IfFileExists{upquote.sty}{\usepackage{upquote}}{}
\IfFileExists{microtype.sty}{% use microtype if available
  \usepackage[]{microtype}
  \UseMicrotypeSet[protrusion]{basicmath} % disable protrusion for tt fonts
}{}
\makeatletter
\@ifundefined{KOMAClassName}{% if non-KOMA class
  \IfFileExists{parskip.sty}{%
    \usepackage{parskip}
  }{% else
    \setlength{\parindent}{0pt}
    \setlength{\parskip}{6pt plus 2pt minus 1pt}}
}{% if KOMA class
  \KOMAoptions{parskip=half}}
\makeatother
\usepackage{longtable,booktabs,array}
\usepackage{calc} % for calculating minipage widths
% Correct order of tables after \paragraph or \subparagraph
\usepackage{etoolbox}
\makeatletter
\patchcmd\longtable{\par}{\if@noskipsec\mbox{}\fi\par}{}{}
\makeatother
% Allow footnotes in longtable head/foot
\IfFileExists{footnotehyper.sty}{\usepackage{footnotehyper}}{\usepackage{footnote}}
\makesavenoteenv{longtable}
\usepackage{graphicx}
\makeatletter
\newsavebox\pandoc@box
\newcommand*\pandocbounded[1]{% scales image to fit in text height/width
  \sbox\pandoc@box{#1}%
  \Gscale@div\@tempa{\textheight}{\dimexpr\ht\pandoc@box+\dp\pandoc@box\relax}%
  \Gscale@div\@tempb{\linewidth}{\wd\pandoc@box}%
  \ifdim\@tempb\p@<\@tempa\p@\let\@tempa\@tempb\fi% select the smaller of both
  \ifdim\@tempa\p@<\p@\scalebox{\@tempa}{\usebox\pandoc@box}%
  \else\usebox{\pandoc@box}%
  \fi%
}
% Set default figure placement to htbp
\def\fps@figure{htbp}
\makeatother
\setlength{\emergencystretch}{3em} % prevent overfull lines
\providecommand{\tightlist}{%
  \setlength{\itemsep}{0pt}\setlength{\parskip}{0pt}}
\usepackage{bookmark}
\IfFileExists{xurl.sty}{\usepackage{xurl}}{} % add URL line breaks if available
\urlstyle{same}
\hypersetup{
  pdftitle={Final report},
  hidelinks,
  pdfcreator={LaTeX via pandoc}}

\title{Final report}
\author{}
\date{\vspace{-2.5em}2026-04-19}

\begin{document}
\maketitle

\section{Abstract}\label{abstract}

\section{Introduction}\label{introduction}

\section{Methods}\label{methods}

\subsubsection{Study design and
participants}\label{study-design-and-participants}

The study follows a cohort study design with data from Electronic health
record (EHR) from a multi-center healthcase system in the United States
including the patient information, comorbidities diagnostic codes, and
residential addresses associated and period of habitation. In addition,
air pollution data from Environmental Protection Agency (EPA) was used
in the study with Air quality index (AQI) values for associated county
and availability of AQI by year in the study.

Participants were included in the cohort if the entry date of the
participant during the index period between January 1st, 2005 and
December 31st, 2010, participants' age of 40 years or older, and
availability of residential data of 10 years minimum prior to the index
date.

\subsubsection{Data handling}\label{data-handling}

The raw dataset underwent a cleaning and transformation process to
ensure clinical interpretability, data validity, and adherence to the
defined inclusion of the criteria. The initial preprocessed cleaning
ensure the standardization of predictors names and data format
consistency; the standardization converted all predictors names to
lowercase convention, and the temporal variables normalized to date
formatting of MM/DD/YYYY. This phase included the reconciliation of
categorical inconsistencies, specifically mixed case entries values in
categorical variables detected in Smoking status variable. In addition,
remapping of binary integers to descriptive strings to facilitate the
interpretation of the analytical results.

Feature engineering was performed for Age variable to transform
continuous variable to appropriate and interpretable bins: \textless40,
40--49, 50--59, 60--69, and 70+ years. Body Mass Index (BMI) was
categorized into four ordinal categories based on World Health
Organization (WHO) definitions.

In order to quantify and assess the environmental risk and exposure, AQI
exposure metric was defined as the weighted air pollution exposure based
on residential address records and EPA monitoring data accounted for
10-year exposure indexed in the study.

Data quality control measures were performed to identify outliers, data
redundancies and missing data. Duplicated observations of participant ID
were removed to ensure unique record for each individual within the
cohort. An outlier in the age value of 345 was identified and recoded as
NA to prevent bias in the models. In order to address the missing data
in the AQI exposure predictor, two approaches were tested: Complete Case
Analysis (CCA) and Multivariate Imputation by Chained Equation (MICE).
The CCA approach was performed to remove the data across all variables
where AQI exposure is missing, where MICE approach was employed to
estimate missing values with statistical estimate. Specifically,
Predictive Mean Matching (PMM) method was employed in MICE approach in
which the imputed values are sampled from observed data points in the
dataset. To account for uncertainty inherent in the imputation process,
we generated 20 independent datasets that include different estimates
for the missing values. For each imputed dataset, a time weighted
average AQI of median AQI value was calculated by integrating the
residential history of each participant.

\subsection{Statistical analysis}\label{statistical-analysis}

The primary association between air pollution exposure and the ten-year
risk of receiving COPD diagnosis was evaluated using Cox Proportional
Hazard models. To ensure the robust of the assessment, the weighted AQI
expsosure predictor was analyzed through three distinct variations: as a
continous numerical scale, data-derived quartiles (Q1-Q4), and according
to regulatory EPA categories (e.g.~``Good'', ``Moderate'', ``Unhealthy
for Sensitive Groups''). All models were adjusted for potential
confounding factors including age, sex, asthma diagnosis, smoking
status.

Statistical evidence was assessed using Hazard Ratios (HR) with
corresponding 95\% confidence intervals (CI), and results visualizaed
with forest plots to compare different exposure variations. Additional
sensitivity analysis was conducted by comparing results from CCA model
against those from MICE imputed model. The resulting estimate from the
MICE imputed model were aggregated using Rubin's Rules; the pooling
procedure reported the point estimates and 95\% CI reflecting the
variance of within imputation and between imputation. This is used to
confirm the stability of the observed risk due to the missing AQI values
of EPA dataset. All statistical tests were two sided and a p-value of
\textless{} 0.05 was considered statistically significant.

\section{Results}\label{results}

\subsubsection{Cohort characteristics}\label{cohort-characteristics}

Following the application of the inclusion criteria, the final
analytical cohort consisted of 1,973 participants, among whom 176 COPD
events were recorded (8.9\% incidence). The baseline demographic and
clinical characteristics of the study population, stratified by COPD
diagnosis, are summarized in Table 1.

The cohort was characterized by a balanced distribution of biological
sex, with males representing 50.1\% of the non-COPD group and 52.3\% of
the COPD group. Analysis of age distribution revealed that the highest
proportions of COPD diagnoses occurred in the 60--69 (42.0\%) and 70--79
(27.3\%) age categories. Clinically, a substantial disparity was
observed in asthma prevalence: 30.1\% of participants in the COPD cohort
had a history of asthma, compared to only 8.8\% in the non-COPD cohort.

Behavioral factors, specifically smoking status, were relatively
comparable across both groups; 60.2\% of COPD cases and 63.7\% of
non-cases reported having ``never'' smoked. However, regarding
environmental exposure, the mean weighted AQI was markedly higher in the
COPD cohort (53.60; SD: 18.69) than in the non-COPD cohort (40.67; SD:
22.55). These descriptive findings suggest that while traditional risk
factors such as sex and smoking status were balanced, age, asthma
history, and environmental pollution levels were notably differentiated
by COPD diagnosis status.

\paragraph{\texorpdfstring{\emph{Table 1 - Cohort
characteristics}}{Table 1 - Cohort characteristics}}\label{table-1---cohort-characteristics}

\begin{longtable}[]{@{}llll@{}}
\toprule\noalign{}
& level & No & Yes \\
\midrule\noalign{}
\endhead
\bottomrule\noalign{}
\endlastfoot
n & & 1797 & 176 \\
Age category (\%) & 40-49 & 207 (11.5) & 3 (1.7) \\
& 50-59 & 542 (30.2) & 32 (18.2) \\
& 60-69 & 683 (38.0) & 74 (42.0) \\
& 70-79 & 309 (17.2) & 48 (27.3) \\
& 80+ & 56 (3.1) & 19 (10.8) \\
Sex (\%) & Female & 897 (49.9) & 84 (47.7) \\
& Male & 900 (50.1) & 92 (52.3) \\
BMI category (\%) & Normal weight & 637 (35.4) & 49 (27.8) \\
& Obesity & 368 (20.5) & 49 (27.8) \\
& Overweight & 661 (36.8) & 70 (39.8) \\
& Underweight & 131 (7.3) & 8 (4.5) \\
Smoking status (\%) & Current & 196 (10.9) & 25 (14.2) \\
& Former & 457 (25.4) & 45 (25.6) \\
& Never & 1144 (63.7) & 106 (60.2) \\
Asthma (\%) & No & 1639 (91.2) & 123 (69.9) \\
& Yes & 158 (8.8) & 53 (30.1) \\
Diabetes (\%) & No & 1568 (87.3) & 143 (81.2) \\
& Yes & 229 (12.7) & 33 (18.8) \\
Hypertension (\%) & No & 1095 (60.9) & 107 (60.8) \\
& Yes & 702 (39.1) & 69 (39.2) \\
AQI Exposure (mean (SD)) & & 40.67 (22.55) & 53.60 (18.69) \\
Socioeconomic status (\%) & Low & 519 (28.9) & 53 (30.1) \\
& Medium & 927 (51.6) & 81 (46.0) \\
& High & 351 (19.5) & 42 (23.9) \\
\end{longtable}

\subsubsection{Primary multivariate analysis: AQI and COPD
risk}\label{primary-multivariate-analysis-aqi-and-copd-risk}

The association between long-term air pollution exposure (based on AQI
Exposure variable) and the COPD risk was evaluated using Cox
Proportional hazards model. After adjusting for demographics variables
(Age, Sex), Clinical related factors(Asthma diagnosis, BMI, Diabetes
status, Hypertension status), Smoking status and socioeconomic status
variables, the weighted AQI exposure remained a significant predictor
for COPD diagnosis with HR: 1.03; 95\% CI: 1.02--1.04; p \textless{}
0.001. This indicates that for every one unit increase in the AQI
exposure, there is a 3\% increase in the risk of COPD diagnosis.

Based on Table 2, the demographic variable such as Age, and clinical
related factors such as Asthma, and BMI showed strong associations
regarding the risk of receiving COPD diagnosis. Asthma was a critical
predictor where a three-fold increase in risk of COPD diagnosis for the
positive diagnosis of Asthma. Regarding age predictor, an increase in
one year of age was associated with 6.2\% increase in risk of receiving
COPD diagnosis, similarly for BMI predictor with HR: 1.06; p \textless{}
0.001. We noted that Smoking status and biological sex predictors were
not statistically significant in the results.

\paragraph{\texorpdfstring{\emph{Table 2 - Primary model
results}}{Table 2 - Primary model results}}\label{table-2---primary-model-results}

\begin{verbatim}
## # A tibble: 11 x 7
##    term                 estimate std.error statistic  p.value conf.low conf.high
##    <chr>                   <dbl>     <dbl>     <dbl>    <dbl>    <dbl>     <dbl>
##  1 aqi_exposure            1.03    0.00387     7.41  1.29e-13    1.02      1.04 
##  2 age                     1.06    0.00832     7.23  4.93e-13    1.04      1.08 
##  3 sexMale                 1.08    0.154       0.527 5.98e- 1    0.802     1.47 
##  4 asthma                  3.48    0.166       7.51  5.94e-14    2.51      4.82 
##  5 smoking_statusFormer    0.903   0.253      -0.405 6.85e- 1    0.550     1.48 
##  6 smoking_statusNever     0.790   0.226      -1.04  2.97e- 1    0.508     1.23 
##  7 bmi                     1.06    0.0153      3.78  1.57e- 4    1.03      1.09 
##  8 ses_statusLow           0.694   0.209      -1.74  8.10e- 2    0.460     1.05 
##  9 ses_statusMedium        0.632   0.192      -2.40  1.64e- 2    0.434     0.919
## 10 diabetes                1.20    0.196       0.923 3.56e- 1    0.816     1.76 
## 11 hypertension            0.803   0.158      -1.39  1.65e- 1    0.590     1.09
\end{verbatim}

\subsubsection{Model diagnostics and
assumptions}\label{model-diagnostics-and-assumptions}

The validity of Cox Proportional Hazards model was confirmed through an
assessment of its underlying assumptions. To evaluate the linearity of
continous covariates, the relationship between the log-hazard of COPD
and continuous predictors---including weighted AQI exposure, BMI, and
age---was examined. Visual inspection of these functional forms
confirmed a linear relationship for all variables, indicating that the
linearity assumption was satisfied without the need for non-linear
transformations. (see Appendix Figure 1)

Furthermore, the proportional hazards assumption was evaluated using
both global and individual Schoenfeld residual diagnostics. he global
Schoenfeld test yielded a non-significant result (p = 0.67), suggesting
that the hazard ratios remained constant over the follow-up period.
Graphical analysis of the Schoenfeld residuals across all predictors
demonstrated no discernible trends or time-dependent effects.
Collectively, these diagnostics confirm that the fundamental assumptions
of the Cox Proportional Hazards model were met, ensuring the reliability
of the reported risk estimates. (see Appendix Figure 2)

\subsubsection{Sensitivity analysis and missing data
reliability}\label{sensitivity-analysis-and-missing-data-reliability}

\textbf{Interaction model}

To evaluate the robustness of the results regarding the impact of air
pollution on COPD risk by sex, an interaction term analysis between AQI
exposure variable and biological sex was performed. Based on the results
in Table 3, the interaction term resulted in HR of 0.9995 with p value
of 0.944640 that signified the interaction term is not statistically
significant. This result suggests that the association between air
pollution exposure and the hazard of receiving a COPD diagnosis does not
differ significantly between males and females in the study population.

\paragraph{\texorpdfstring{\emph{Table 3 - Interaction term AQI exposure
and Sex
model}}{Table 3 - Interaction term AQI exposure and Sex model}}\label{table-3---interaction-term-aqi-exposure-and-sex-model}

\begin{verbatim}
## # A tibble: 12 x 7
##    term                 estimate std.error statistic  p.value conf.low conf.high
##    <chr>                   <dbl>     <dbl>     <dbl>    <dbl>    <dbl>     <dbl>
##  1 aqi_exposure            1.03    0.00575    5.03   4.93e- 7    1.02      1.04 
##  2 age                     1.06    0.00833    7.23   4.99e-13    1.04      1.08 
##  3 sexMale                 1.12    0.440      0.249  8.03e- 1    0.471     2.64 
##  4 asthma                  3.48    0.166      7.51   6.04e-14    2.51      4.82 
##  5 smoking_statusFormer    0.904   0.254     -0.397  6.91e- 1    0.550     1.49 
##  6 smoking_statusNever     0.792   0.227     -1.03   3.03e- 1    0.508     1.23 
##  7 bmi                     1.06    0.0153     3.78   1.57e- 4    1.03      1.09 
##  8 ses_statusLow           0.694   0.210     -1.74   8.17e- 2    0.460     1.05 
##  9 ses_statusMedium        0.632   0.192     -2.40   1.64e- 2    0.434     0.919
## 10 diabetes                1.20    0.196      0.924  3.56e- 1    0.816     1.76 
## 11 hypertension            0.803   0.158     -1.39   1.65e- 1    0.590     1.09 
## 12 aqi_exposure:sexMale    0.999   0.00775   -0.0694 9.45e- 1    0.984     1.01
\end{verbatim}

\textbf{Numerical, quartile and EPA AQI exposure predictor}

The sensitivity of the findings was further explored by modeling the AQI
exposure in three different definitions: Continous numerical scale, data
derived quartiles, and EPA recommended categories. Based on Figure 1, in
the continuous model, every one unit increase in AQI was associated with
2.8\% increase in risk of receiving COPD diagnosis. The quartile
analysis improved the interpretation of the results by revealing that
the individuals in highest AQI exposure group (Q4) faced 5.2 times the
hazard compared to those in the lowest group (Q1). Similarly, using EPA
categories, ``Moderate'' air pollution exposure increased the risk by
2.78 times compared to Good air quality, while ``Unhealthy for Sensitive
Groups'' showed a much high hazard ratio of 6.3. However, the category
``Unhealthy for Sensitive Groups'' is noted to not be statistically
significant with p value of 0.07.

\paragraph{\texorpdfstring{\emph{Figure 1 - Comparison between
Numerical, quartile and EPA AQI exposure predictor in
models}}{Figure 1 - Comparison between Numerical, quartile and EPA AQI exposure predictor in models}}\label{figure-1---comparison-between-numerical-quartile-and-epa-aqi-exposure-predictor-in-models}

\begin{verbatim}
## # A tibble: 6 x 8
##   Model_Type term         HR Lower_95 Upper_95 Model_Label term_clean label_text
##   <chr>      <chr>     <dbl>    <dbl>    <dbl> <chr>       <chr>      <chr>     
## 1 1          aqi_cate~  2.79    2.02      3.85 EPA Catego~ Moderate   2.79 (2.0~
## 2 1          aqi_cate~  6.32    0.856    46.7  EPA Catego~ Unhealthy~ 6.32 (0.8~
## 3 2          aqi_quar~  1.93    1.04      3.59 AQI Quarti~ Q2         1.93 (1.0~
## 4 2          aqi_quar~  3.57    2.03      6.26 AQI Quarti~ Q3         3.57 (2.0~
## 5 2          aqi_quar~  5.36    3.11      9.23 AQI Quarti~ Q4 (Highe~ 5.36 (3.1~
## 6 3          aqi_expo~  1.03    1.02      1.04 Continuous~ aqi_expos~ 1.03 (1.0~
\end{verbatim}

\pandocbounded{\includegraphics[keepaspectratio]{Final-report_files/figure-latex/model analysis Comparaison-1.pdf}}

\textbf{CCA vs MICE}

Finally, reliability of the analysis is assessed regarding the missing
data treatment by comapring the Complete Case Analysis (CCA) with
Multivariate Imputation by Chained Equation (MICE). Based on Figure 2,
the results from both approaches are similar with both models indicating
that a one unit increase in AQI exposure increases the risk of receiving
COPD diagnosis by 3-4\%. Since the difference between the CCA and MICE
is minimal, the model results are reliable and the missing data in
weighted AQI exposure is likely Missing completely at random or missing
at random.

\paragraph{\texorpdfstring{\emph{Figure 2 - Comparison between MICE and
CCA methods of treating missing
data}}{Figure 2 - Comparison between MICE and CCA methods of treating missing data}}\label{figure-2---comparison-between-mice-and-cca-methods-of-treating-missing-data}

\pandocbounded{\includegraphics[keepaspectratio]{Final-report_files/figure-latex/model analysis CCA vs MICE vs primary-1.pdf}}

\section{Discussion}\label{discussion}

\section{References}\label{references}

\section{Appendix}\label{appendix}

\pandocbounded{\includegraphics[keepaspectratio]{Final-report_files/figure-latex/model assumption-1.pdf}}
\pandocbounded{\includegraphics[keepaspectratio]{Final-report_files/figure-latex/model assumption-2.pdf}}

\end{document}
